\subsection{The first unresolved wild coefficient}
The next theorem narrows the first coefficient beyond the exact wild bands to one necessary configuration of actual torsion representations.

\begin{theorem}[The first coefficient beyond the exact bands]
\label{thm:first-boundary-normal-form}
Let $p\ge5$ and $r\ge1$, and put
\[
 k_0=\frac{p^2+1}{2},\qquad b_0=\frac{p-1}{2},\qquad
 e=k_0(p-1)p^r,\qquad n=r+3.
\]
Suppose that two groups in the standing height-two filtration class
over a field of absolute ramification index $e$ have isomorphic
actual $p^n$-torsion and different special inertia characters.
Then both are whole potentially supersingular groups with two-slope
good models. In the common tame normalization their Hasse parameters
and special tangent scalars satisfy
\[
 z_1=z_2=b_0p^n,\qquad \frac{\alpha_2}{\alpha_1}=-1.
\]
Both residual $\mathbf F_p$-representations over $K$ are split.
Their canonical and complementary lines have distinct inertia
characters, and the given isomorphism sends each canonical line
to the complementary line of the other representation.
\end{theorem}

This is a necessary form for an assumed unequal-character comparison.
It does not assert that this form occurs or determine the character
index at $e=k_0(p-1)p^r$.

Section~\ref{sec:first-boundary-normal-form} proves the result using the general two-chain congruence and the actual quotient markings.

\begin{corollary}[Point fields at the first boundary]
\label{cor:boundary-point-field-degrees}
Under the assumptions of
Theorem~\ref{thm:first-boundary-normal-form}, pass to the completed
maximal unramified base and the common tame good field $L$.
For every exact $p^j$-point $Q$ of either group whose $p^{j-1}$-multiple is a
nonzero canonical point, write $[L(Q):L]=p^d$. Then
\[
 d\ge\max\{0,2j-1-r\}\quad(2\le j<n),\qquad
 d\ge r+5\quad(j=n).
\]
Each relative step at levels $r+1,r+2,r+3$ has degree $p^2$.
In particular every canonical $p^{n-1}$-point field has degree
at least $p^{r+3}$. If $r=1$ and $n=4$, the degrees at levels
$p^2,p^3,p^4$ are respectively
\[
 p^2,\qquad p^4,\qquad p^6.
\]
\end{corollary}

Proposition~\ref{prop:proper-canonical-point-fields} proves a
more general proper-level bound for an exchanged central
comparison. The point-field restrictions above do not settle
whether the boundary comparison occurs.

\begin{corollary}[Nonsplit Cartan exclusion at the first boundary]
\label{cor:boundary-nonsplit-cartan-image}
Let $p\ge5$, $r\ge1$, $n=r+3$, and
\[
                    e=\frac{p^2+1}{2}(p-1)p^r.
\]
Suppose that two groups in the standing height-two filtration class
over a field of absolute index $e$ have isomorphic actual
$p^n$-torsion and different special inertia characters.
Pass to the completed maximal unramified base and its common
tame good extension $L$ of degree $p^2-1$.
For each group, the finite image of
$\operatorname{Gal}(\overline L/L)$ on its $p^n$-torsion is not
contained in any conjugate of
\[
                    (1+p\mathcal O)/(1+p^n\mathcal O),
\]
where $\mathcal O$ is the unramified quadratic
$\mathbf Z_p$-algebra and $\mathcal O/p^n\mathcal O$ acts on
$(\mathbf Z/p^n\mathbf Z)^2$ through its regular representation.
\end{corollary}

This is a necessary restriction on an assumed actual comparison.
It does not exclude the remaining non-Cartan images or determine
the boundary character index. Section~\ref{sec:boundary-nonsplit-image-proof}
proves the restriction by constructing an actual invariant-ring map
under the additional image-containment assumption.

\begin{corollary}[The irreducible quadratic self-twist case]
\label{cor:boundary-nonsplit-self-twist}
Under the assumptions of
Corollary~\ref{cor:boundary-nonsplit-cartan-image}, the finite
representation of either group admits no tame quadratic self-twist
with an invertible intertwiner whose reduction has irreducible
quadratic characteristic polynomial over $\mathbf F_p$.
\end{corollary}


The proof in Section~\ref{sec:boundary-nonsplit-image-proof} covers
arbitrary perfect residue fields. The split-intertwiner case is not
part of this corollary.
